\documentclass[11pt, a4paper]{article}

\usepackage[margin=2.5cm]{geometry}
\usepackage{booktabs}
\usepackage{array}
\usepackage{listings}
\usepackage{xcolor}
\usepackage{hyperref}
\usepackage{parskip}
\usepackage{titlesec}
\usepackage{enumitem}
\usepackage{caption}

\hypersetup{
    colorlinks=true,
    linkcolor=blue,
    urlcolor=blue,
}

\definecolor{jsonbg}{RGB}{245,245,245}
\definecolor{jsonkw}{RGB}{0,0,180}
\definecolor{jsonstr}{RGB}{180,0,0}
\definecolor{jsoncmt}{RGB}{80,120,80}

\lstdefinelanguage{json}{
    basicstyle=\ttfamily\small,
    backgroundcolor=\color{jsonbg},
    frame=single,
    rulecolor=\color{gray!40},
    numbers=none,
    breaklines=true,
    showstringspaces=false,
    literate=
        *{0}{{{\color{jsonkw}0}}}{1}
         {1}{{{\color{jsonkw}1}}}{1}
         {2}{{{\color{jsonkw}2}}}{1}
         {3}{{{\color{jsonkw}3}}}{1}
         {4}{{{\color{jsonkw}4}}}{1}
         {5}{{{\color{jsonkw}5}}}{1}
         {6}{{{\color{jsonkw}6}}}{1}
         {7}{{{\color{jsonkw}7}}}{1}
         {8}{{{\color{jsonkw}8}}}{1}
         {9}{{{\color{jsonkw}9}}}{1}
         {:}{{{\color{black}{:}}}}{1}
         {,}{{{\color{black}{,}}}}{1}
         {\{}{{{\color{black}{\{}}}}{1}
         {\}}{{{\color{black}{\}}}}}{1}
         {[}{{{\color{black}{[}}}}{1}
         {]}{{{\color{black}{]}}}}{1},
}

\titleformat{\section}{\large\bfseries}{}{0em}{}[\titlerule]
\titleformat{\subsection}{\normalsize\bfseries}{}{0em}{}

\title{\textbf{Oscilloscope Benchmark Guide}\\
\large Automated Behavioural Training Rig --- Real-Time Performance Characterisation}
\author{Joel Kiener}
\date{Bachelor Thesis, FS 2026}

\begin{document}

\maketitle

\tableofcontents
\newpage

% ==============================================================================
\section*{GPIO Quick Reference}
\addcontentsline{toc}{section}{GPIO Quick Reference}
% ==============================================================================

All GPIO numbers are BCM. Connect the oscilloscope GND probe to any Pi GND
pin (e.g.\ physical pin 6 or 9). Logic level is 3.3\,V on all GPIO.
Beams are \textbf{active low} (beam broken $\rightarrow$ pin goes LOW).
LEDs and valves are \textbf{active high} (on $\rightarrow$ pin goes HIGH).

\begin{table}[h!]
\centering
\begin{tabular}{llll}
\toprule
\textbf{Signal} & \textbf{BCM} & \textbf{Physical} & \textbf{Logic} \\
\midrule
Beam left    & 2  & 3  & Active low  \\
Beam center  & 3  & 5  & Active low  \\
Beam right   & 4  & 7  & Active low  \\
LED left     & 13 & 33 & Active high \\
LED center   & 19 & 35 & Active high \\
LED right    & 26 & 37 & Active high \\
Valve left   & 0  & 27 & Active high \\
Valve right  & 5  & 29 & Active high \\
Audio L/R    & \multicolumn{3}{l}{3.5\,mm jack (tip = left, ring = right)} \\
\bottomrule
\end{tabular}
\caption{GPIO pin reference (\texttt{RPi\_main/config.py}).}
\end{table}

\newpage

% ==============================================================================
\section{FSM Transition Latency (Headline Number)}
% ==============================================================================

\subsection*{What it measures}

The end-to-end time from a hardware beam break to the Pi asserting a GPIO
output in response. This captures the full software chain:

\[
\text{kernel IRQ} \;\rightarrow\; \text{Python callback} \;\rightarrow\;
\text{event queue} \;\rightarrow\; \text{FSM thread} \;\rightarrow\; \text{GPIO write}
\]

This is the number most directly comparable to pyControl
($556 \pm 17\,\mu\text{s}$, Akam et al.\ 2022) and LabNet
($0.26\,\text{ms}$ median, Schatz \& Winter 2022).

\subsection*{Hardware setup}

\begin{itemize}[noitemsep]
    \item \textbf{CH1}: Beam center (BCM 3, physical 5)
    \item \textbf{CH2}: LED center (BCM 19, physical 35)
    \item \textbf{GND}: Pi physical pin 6
    \item To simulate a beam break cleanly: connect a short wire from the beam
          pin to GND, or use a signal generator at 3.3\,V logic. Alternatively,
          block the IR beam with a card.
\end{itemize}

\subsection*{Trial definition}

Create a new substage called \texttt{bench\_fsm} in the curriculum builder and
paste the following task config:

\begin{lstlisting}[language=json]
{
  "trial_id": "bench_fsm_latency",
  "side_mode": "fixed_left",
  "base_iti_s": 1.0,
  "fail_iti_s": 1.0,
  "initial_state": "wait_for_beam",
  "states": [
    {
      "id": "wait_for_beam",
      "duration": 30.0,
      "entry_actions": [],
      "exit_actions": [],
      "transitions": [
        { "trigger": "beam_break", "target": "center", "next_state": "respond" },
        { "trigger": "timeout", "next_state": "__wrong__" }
      ]
    },
    {
      "id": "respond",
      "duration": 2.0,
      "entry_actions": [ { "type": "led_on",  "target": "center" } ],
      "exit_actions":  [ { "type": "led_off", "target": "center" } ],
      "transitions": [
        { "trigger": "timeout", "next_state": "__correct__" }
      ]
    }
  ]
}
\end{lstlisting}

\subsection*{Oscilloscope settings}

\begin{itemize}[noitemsep]
    \item Timebase: 1\,ms/div
    \item CH1: DC coupling, 2\,V/div --- beam signal falls from 3.3\,V to 0\,V
    \item CH2: DC coupling, 2\,V/div --- LED rises from 0\,V to 3.3\,V
    \item Trigger: CH1, falling edge, $\sim$1.5\,V threshold
    \item Mode: Normal (single-shot per break event)
    \item Use cursor $\Delta t$ or built-in measurement to read CH1 falling
          $\rightarrow$ CH2 rising
\end{itemize}

\subsection*{Procedure}

Start a session, assign a cage to substage \texttt{bench\_fsm}, and start
continuous trials. Arm the scope in Normal mode, break the beam, read
$\Delta t$, record the value, reset, repeat. Collect $N \geq 500$ samples.
Most bench scopes (Rigol, Siglent, Keysight) have a Statistics mode under
Measure that accumulates mean/std/min/max automatically.

\subsection*{What to report}

Mean $\pm$ std, minimum, maximum, 99th percentile, histogram (100\,$\mu$s bins).
Repeat under 12-cage load (Section~\ref{sec:load}) and compare.

\newpage

% ==============================================================================
\section{Threading.Timer Accuracy (State Timeout Jitter)}
% ==============================================================================

\subsection*{What it measures}

FSM state timeouts use Python's \texttt{threading.Timer}. On a non-RTOS Linux
system, timer callbacks are subject to OS scheduler jitter. This measures how
accurately a programmed timeout actually fires, which affects stimulus duration
precision and spike-event alignment.

\subsection*{Hardware setup}

\begin{itemize}[noitemsep]
    \item \textbf{CH1}: LED center (BCM 19, physical 35)
    \item LED turns ON at state entry, OFF when timeout fires
    \item No beam break needed
\end{itemize}

\subsection*{Trial definition}

\begin{lstlisting}[language=json]
{
  "trial_id": "bench_timeout",
  "side_mode": "fixed_left",
  "base_iti_s": 0.5,
  "fail_iti_s": 0.5,
  "initial_state": "timed_state",
  "states": [
    {
      "id": "timed_state",
      "duration": 1.0,
      "entry_actions": [ { "type": "led_on",  "target": "center" } ],
      "exit_actions":  [ { "type": "led_off", "target": "center" } ],
      "transitions": [
        { "trigger": "timeout", "next_state": "__correct__" }
      ]
    }
  ]
}
\end{lstlisting}

Test with \texttt{duration} = 0.1\,s, 0.5\,s, 1.0\,s, and 2.0\,s.
Shorter timeouts will show higher relative jitter.

\subsection*{Oscilloscope settings}

\begin{itemize}[noitemsep]
    \item Timebase: 200\,ms/div for 1\,s duration (2\,s total window)
    \item CH1: DC coupling, 2\,V/div
    \item Trigger: CH1 rising edge
    \item Use automatic pulse-width measurement in Statistics mode
\end{itemize}

\subsection*{What to report}

Mean measured duration vs.\ programmed duration (systematic offset), standard
deviation (jitter), min/max. Table across all four duration values.
Repeat under load (Section~\ref{sec:load}).

\newpage

% ==============================================================================
\section{Hold Timer Accuracy}
% ==============================================================================

\subsection*{What it measures}

The \texttt{hold\_ms} feature requires the animal to maintain a beam break for
a specified duration before the FSM advances. This timer is started at
beam-break time. Its accuracy directly affects task validity: if
\texttt{hold\_ms=100} fires at 85\,ms, animals are rewarded for shorter holds
than intended.

\subsection*{Hardware setup}

\begin{itemize}[noitemsep]
    \item \textbf{CH1}: Beam center (BCM 3, physical 5)
    \item \textbf{CH2}: LED center (BCM 19, physical 35) --- turns on after hold completes
    \item Simulate a held break with a jumper wire from BCM 3 to GND (do not
          use a finger, which may wobble and trigger debounce events)
\end{itemize}

\subsection*{Trial definition}

\begin{lstlisting}[language=json]
{
  "trial_id": "bench_hold",
  "side_mode": "fixed_left",
  "base_iti_s": 0.5,
  "fail_iti_s": 0.5,
  "initial_state": "wait",
  "states": [
    {
      "id": "wait",
      "duration": 30.0,
      "entry_actions": [],
      "exit_actions": [],
      "transitions": [
        {
          "trigger": "beam_break",
          "target": "center",
          "hold_ms": 100,
          "next_state": "held"
        },
        { "trigger": "timeout", "next_state": "__wrong__" }
      ]
    },
    {
      "id": "held",
      "duration": 1.0,
      "entry_actions": [ { "type": "led_on",  "target": "center" } ],
      "exit_actions":  [ { "type": "led_off", "target": "center" } ],
      "transitions": [
        { "trigger": "timeout", "next_state": "__correct__" }
      ]
    }
  ]
}
\end{lstlisting}

Test with \texttt{hold\_ms} = 50, 100, 200, 500.

\subsection*{Oscilloscope settings}

\begin{itemize}[noitemsep]
    \item Timebase: 10\,ms/div for 100\,ms hold
    \item Trigger: CH1 falling edge
    \item Measure $\Delta t$: CH1 falling edge $\rightarrow$ CH2 rising edge
\end{itemize}

\subsection*{What to report}

Mean hold duration vs.\ programmed \texttt{hold\_ms} (systematic bias),
std deviation, table across all four values. $N \geq 300$ per value.

\newpage

% ==============================================================================
\section{Audio Onset Latency}
% ==============================================================================

\subsection*{What it measures}

The time from the FSM dispatching a \texttt{play\_clicks} action (LED GPIO as
proxy timestamp) to the first audio sample appearing at the DAC output
(3.5\,mm jack). This is the \texttt{sounddevice} buffer latency.

This is the most experiment-critical latency: the animal hears the first click
some time after state entry. All trial timestamps are aligned to state entry,
so audio onset latency is a systematic offset in all behavioural data that
must be reported so others can correct for it.

\subsection*{Hardware setup}

\begin{itemize}[noitemsep]
    \item \textbf{CH1}: LED left (BCM 13, physical 33) --- state entry marker
    \item \textbf{CH2}: 3.5\,mm jack tip (left channel)
          Use a 3.5\,mm-to-BNC adapter, or strip a cable and probe the tip wire.
    \item CH2: AC coupling, 200\,mV/div (audio is $\pm$1\,V AC waveform)
\end{itemize}

\subsection*{Trial definition}

\begin{lstlisting}[language=json]
{
  "trial_id": "bench_audio_latency",
  "side_mode": "fixed_left",
  "base_iti_s": 1.0,
  "fail_iti_s": 1.0,
  "initial_state": "play",
  "states": [
    {
      "id": "play",
      "duration": 3.0,
      "entry_actions": [
        { "type": "led_on",      "target": "left" },
        { "type": "play_clicks", "left_rate": 100, "right_rate": 0,
          "click_duration": 2.5 }
      ],
      "exit_actions": [
        { "type": "led_off", "target": "left" }
      ],
      "transitions": [
        { "trigger": "timeout", "next_state": "__correct__" }
      ]
    }
  ]
}
\end{lstlisting}

\texttt{left\_rate=100} gives dense clicks that are easy to identify on the
scope. \texttt{right\_rate=0} isolates the left channel for clean probing.

\subsection*{Oscilloscope settings}

\begin{itemize}[noitemsep]
    \item Timebase: 5\,ms/div (50\,ms window; audio latency is typically 10--30\,ms)
    \item CH1: DC coupling, 2\,V/div
    \item CH2: AC coupling, 500\,mV/div
    \item Trigger: CH1 rising edge
    \item Measure $\Delta t$: CH1 rising edge $\rightarrow$ first audio peak on CH2
\end{itemize}

\subsection*{What to report}

Mean audio onset latency $\pm$ std ($N \geq 200$). State this value explicitly
in the thesis Methods section as the systematic offset between logged state
entry time and actual stimulus delivery.

\newpage

% ==============================================================================
\section{Click Inter-Arrival Interval Distribution}
% ==============================================================================

\subsection*{What it measures}

Whether the intended Poisson-distributed inter-click intervals (ICI) are
preserved through the audio pipeline. If \texttt{sounddevice} introduces
timing distortion, click trains are not truly Poisson, which would invalidate
the psychophysics task.

\subsection*{Hardware setup}

\begin{itemize}[noitemsep]
    \item \textbf{CH1}: 3.5\,mm jack tip (left channel), AC coupling
    \item No GPIO probe needed
\end{itemize}

\subsection*{Procedure}

Capture a full 2.5\,s click train waveform at the scope's maximum sample rate.
Export the waveform as CSV via the scope's USB port (supported on most Rigol/
Siglent models). Then in Python:

\begin{enumerate}[noitemsep]
    \item Threshold the waveform to find click peak times.
    \item Compute inter-click intervals (ICI).
    \item Plot ICI histogram and overlay the expected exponential distribution
          with rate $\lambda = $ \texttt{left\_rate}.
    \item Run a Kolmogorov--Smirnov test against the exponential null hypothesis.
\end{enumerate}

Repeat for rate pairs: (80, 20), (60, 40), (50, 50), (100, 0).

\subsection*{What to report}

ICI histogram vs.\ theoretical exponential for each rate, KS test $p$-value,
and any systematic deviation such as the minimum ICI floor imposed by
\texttt{CLICK\_WIDTH\_S = 3\,ms} in \texttt{config.py}.

\newpage

% ==============================================================================
\section{Valve Open Duration Precision}
% ==============================================================================

\subsection*{What it measures}

Reward valves are opened for a fixed duration to deliver a precise water
volume. Duration jitter translates directly to water volume variance, affecting
both welfare records and experiment reproducibility.

\subsection*{Hardware setup}

\begin{itemize}[noitemsep]
    \item \textbf{CH1}: Valve left (BCM 0, physical 27)
\end{itemize}

\subsection*{Trial definition}

\begin{lstlisting}[language=json]
{
  "trial_id": "bench_valve",
  "side_mode": "fixed_left",
  "base_iti_s": 2.0,
  "fail_iti_s": 2.0,
  "initial_state": "reward",
  "states": [
    {
      "id": "reward",
      "duration": 0.15,
      "entry_actions": [ { "type": "valve_open",  "target": "left" } ],
      "exit_actions":  [ { "type": "valve_close", "target": "left" } ],
      "transitions": [
        { "trigger": "timeout", "next_state": "__correct__" }
      ]
    }
  ]
}
\end{lstlisting}

\texttt{duration = 0.15\,s} matches \texttt{VALVE\_OPEN\_DEFAULT\_MS = 150\,ms}
in \texttt{config.py}. Also test with 0.05\,s and 0.30\,s to check linearity.

\subsection*{Oscilloscope settings}

\begin{itemize}[noitemsep]
    \item Timebase: 50\,ms/div
    \item CH1: DC coupling, 2\,V/div
    \item Trigger: CH1 rising edge
    \item Use automatic pulse-width measurement in Statistics mode
\end{itemize}

\subsection*{What to report}

Mean valve open duration vs.\ programmed duration, std deviation, $N \geq 300$.
Convert jitter to volume: a typical solenoid valve delivers approximately
3--8\,$\mu$L/ms, so 1\,ms jitter corresponds to 3--8\,$\mu$L volume variance
per reward.

\newpage

% ==============================================================================
\section{Audio--LED Synchronisation (Multi-modal Stimulus Onset)}
% ==============================================================================

\subsection*{What it measures}

In the actual task, LED and click train are presented simultaneously as the
stimulus. This measures how well-synchronised they are at the hardware output
level, i.e.\ the audio-visual asynchrony experienced by the animal.

\subsection*{Hardware setup}

\begin{itemize}[noitemsep]
    \item \textbf{CH1}: LED left (BCM 13, physical 33)
    \item \textbf{CH2}: 3.5\,mm jack tip (left channel), AC coupling
\end{itemize}

Use the same trial definition as Section~4 (\texttt{bench\_audio\_latency}).
The delta between CH1 rising edge and the first audio click on CH2 is the
audio-visual stimulus onset asynchrony (SOA).

\subsection*{What to report}

Mean SOA $\pm$ std. State this as: ``visual and auditory stimuli are offset by
$X$\,ms, with the LED preceding the first click,'' and note that the animal
perceives the clicks with this delay relative to the logged state entry
timestamp.

\newpage

% ==============================================================================
\section{Camera--GPIO Timestamp Synchronisation}
% ==============================================================================

\subsection*{What it measures}

The Pi embeds a GPIO state snapshot in every UDP frame header, timestamped via
\texttt{CLOCK\_MONOTONIC}. This validates whether the camera frame timestamps
correctly correspond to the actual moment of GPIO state change.

When an LED turns on at time $T_\text{GPIO}$, the first camera frame in which
the LED appears lit should have a header timestamp satisfying:

\[
0 \leq T_\text{frame} - T_\text{GPIO} \leq \frac{1}{60\,\text{fps}} = 16.7\,\text{ms}
\]

Errors outside this range indicate a bug in the timestamp anchoring code in
\texttt{RPi\_main/streamer.py}.

\subsection*{Hardware setup}

\begin{itemize}[noitemsep]
    \item \textbf{CH1}: LED center (BCM 19, physical 35)
    \item Point the cage camera directly at the LED so it is clearly visible
          in-frame.
    \item No second scope probe needed --- the camera is the second channel.
\end{itemize}

\subsection*{Procedure}

\begin{enumerate}[noitemsep]
    \item Run the \texttt{bench\_fsm\_latency} trial (Section~1) with camera
          recording enabled (NAS write on).
    \item Add a single log line in \texttt{engine.py} at the GPIO write call to
          record \texttt{clock\_gettime(CLOCK\_MONOTONIC)} as $T_\text{GPIO}$.
    \item After the session, open the \texttt{.bin} recording in
          \texttt{bin\_viewer.py}.
    \item Find the frame where the LED state bit first changes to 1 in the header.
          Read the embedded timestamp $T_\text{frame}$.
    \item Compute $\Delta t = T_\text{frame} - T_\text{GPIO}$.
    \item Repeat for $N \geq 100$ LED events.
\end{enumerate}

\subsection*{What to report}

Mean synchronisation error $\pm$ std, histogram. Confirm that errors are always
positive (frame timestamp always after GPIO event). Negative values indicate a
timestamp anchoring error.

\newpage

% ==============================================================================
\section{Latency Under Load (Scaling Validation)}
\label{sec:load}
% ==============================================================================

\subsection*{What it measures}

All measurements above should be repeated under full 12-cage load to verify
that the system scales without performance degradation. This is the key
scalability claim of the thesis.

\subsection*{Setup}

\begin{enumerate}[noitemsep]
    \item Run Sections 1, 2, and 4 with only 1 cage active (baseline).
    \item Repeat with 12 cages running simultaneous continuous trials.
          The measured cage runs the benchmark trial definition; the other 11
          run the normal task or any continuous trial loop.
\end{enumerate}

Also check the acquisition logs for UDP sequence number gaps during the
12-cage run (logged by \texttt{watchdog.py}).

\subsection*{What to report}

Table of all metrics at $N=1$ vs.\ $N=12$ cages:

\begin{table}[h!]
\centering
\begin{tabular}{lcccc}
\toprule
\textbf{Metric} & \multicolumn{2}{c}{\textbf{1 cage}} & \multicolumn{2}{c}{\textbf{12 cages}} \\
                & mean & 99th pct & mean & 99th pct \\
\midrule
FSM transition latency     & & & & \\
Timeout jitter (std)       & & & & \\
Audio onset latency        & & & & \\
UDP frame drop rate        & & & & \\
\bottomrule
\end{tabular}
\caption{Real-time performance metrics under single-cage and full 12-cage load.}
\end{table}

\newpage

% ==============================================================================
\section{TCP Command Latency (PC $\rightarrow$ Pi Round-Trip)}
% ==============================================================================

\subsection*{What it measures}

The time from the PC sending a trial start command over TCP to the Pi's first
GPIO output. This captures: network round-trip time + TCP stack latency +
Python socket receive + FSM initialisation.

\subsection*{Hardware setup}

\begin{itemize}[noitemsep]
    \item \textbf{CH1}: LED center (BCM 19, physical 35)
    \item On the PC side: add a \texttt{time.monotonic()} log in
          \texttt{command/tcp\_command\_sender.py} at the moment bytes are
          written to the socket.
    \item Compare the PC-side send timestamp to the scope CH1 rising edge
          timestamp using the scope's time display.
\end{itemize}

Use any trial definition where the initial state has an immediate LED entry
action (e.g.\ \texttt{bench\_fsm\_latency}).

\subsection*{What to report}

Mean PC$\rightarrow$Pi GPIO latency $\pm$ std ($N \geq 200$). Cross-check with
network RTT alone: \texttt{ping -c 1000 192.168.1.101} gives the pure network
component. The difference is the Python processing overhead.

\newpage

% ==============================================================================
\section{SCHED\_FIFO Effect on Jitter}
% ==============================================================================

\subsection*{What it measures}

Whether enabling POSIX real-time scheduling (\texttt{SCHED\_FIFO}) in
\texttt{engine.py} meaningfully reduces latency or jitter. This requires
running \texttt{main.py} as root on the Pi (\texttt{sudo python main.py}).

\subsection*{Procedure}

\begin{enumerate}[noitemsep]
    \item Run Section~1 (FSM transition latency) in \textbf{normal scheduling}
          mode. Record full distribution ($N \geq 500$).
    \item Enable \texttt{SCHED\_FIFO} for the FSM thread in \texttt{engine.py}
          and restart with \texttt{sudo}.
    \item Run Section~1 again under identical conditions.
    \item Compare distributions, especially the tail (99th percentile and max).
\end{enumerate}

\subsection*{What to report}

Side-by-side comparison: does the \emph{mean} change? (Probably not.)
Does the 99th percentile or maximum change? (Likely yes --- \texttt{SCHED\_FIFO}
primarily reduces outlier latencies, not the mean.) Provide a practical
recommendation: is \texttt{SCHED\_FIFO} necessary for this application?

\newpage

% ==============================================================================
\section{Inter-Trial Interval Precision}
% ==============================================================================

\subsection*{What it measures}

The ITI is implemented as \texttt{time.sleep()} in \texttt{ui/cage\_runner.py}
on the PC, followed by a new TCP command to the Pi. This measures how
accurately the ITI executes end-to-end, including network overhead.

\subsection*{Hardware setup}

\begin{itemize}[noitemsep]
    \item \textbf{CH1}: LED center (BCM 19, physical 35)
    \item LED turns on at each new trial's initial state entry.
    \item Measure: time from LED falling edge (trial end) to next LED rising
          edge (next trial start) = actual ITI.
\end{itemize}

Use \texttt{bench\_fsm\_latency} with \texttt{base\_iti\_s = 1.0} and a very
short trial (beam break immediately ends the trial).

\subsection*{Oscilloscope settings}

\begin{itemize}[noitemsep]
    \item Timebase: 500\,ms/div (for 1\,s ITI --- gives $\sim$5\,s window)
    \item Use automatic period or pulse-separation measurement
\end{itemize}

\subsection*{What to report}

Mean ITI vs.\ programmed ITI, std deviation. This sets the floor on how
precisely session pacing can be controlled.

\newpage

% ==============================================================================
\section*{Summary Table}
\addcontentsline{toc}{section}{Summary Table}
% ==============================================================================

\begin{table}[h!]
\centering
\small
\begin{tabular}{clllr}
\toprule
\textbf{\#} & \textbf{Measurement} & \textbf{CH1} & \textbf{CH2} & \textbf{N} \\
\midrule
1  & FSM transition latency        & Beam center (BCM 3)   & LED center (BCM 19) & 500+ \\
2  & Timeout jitter                & LED center (BCM 19)   & ---                 & 500+ \\
3  & Hold timer accuracy           & Beam center (BCM 3)   & LED center (BCM 19) & 300+ \\
4  & Audio onset latency           & LED left (BCM 13)     & Audio jack          & 200+ \\
5  & Click ICI distribution        & Audio jack            & ---                 & 1 rec. \\
6  & Valve duration precision      & Valve left (BCM 0)    & ---                 & 300+ \\
7  & Audio--LED synchronisation    & LED left (BCM 13)     & Audio jack          & 200+ \\
8  & Camera--GPIO sync             & LED center (BCM 19)   & (software)          & 100+ \\
9  & Latency under load            & (repeat 1, 2, 4)      & ---                 & 500+ \\
10 & TCP command latency           & LED center (BCM 19)   & (PC log)            & 200+ \\
11 & SCHED\_FIFO effect            & Beam center (BCM 3)   & LED center (BCM 19) & 500+$\times$2 \\
12 & ITI precision                 & LED center (BCM 19)   & ---                 & 200+ \\
\bottomrule
\end{tabular}
\caption{All measurements at a glance.}
\end{table}

\newpage

% ==============================================================================
\section*{Practical Tips}
\addcontentsline{toc}{section}{Practical Tips}
% ==============================================================================

\subsection*{Probe connections}

Pi GPIO runs at 3.3\,V logic. Use 10$\times$ probe attenuation. Keep the
oscilloscope GND lead short (10\,cm) connected to a Pi GND pin; long ground
leads introduce noise at this voltage level. The beam sensor is active-low
with an internal pull-up, so you will see a clean TTL falling edge on beam
break --- fine for ms-scale measurements.

\subsection*{Collecting statistics automatically}

Most bench scopes (Rigol DS1054Z, Siglent SDS1104X-E, Keysight DSOX1204G)
have a Statistics mode under \textbf{Measure $\rightarrow$ Statistics}. Enable
it, set it to accumulate, and walk away for 10 minutes to collect 500+ samples
automatically. Photograph the statistics screen, or export data via USB as CSV
for offline analysis.

\subsection*{Triggering on beam breaks}

The beam sensor is active-low with 1\,ms software debounce. Set scope trigger
to falling edge, $\sim$1.5\,V threshold to capture the first edge cleanly.

\subsection*{Recommended measurement order}

The following order minimises probe re-wiring and can be completed in one
afternoon session:

\begin{enumerate}[noitemsep]
    \item \textbf{Section 1} (FSM latency) --- main result, beam + LED probes.
    \item \textbf{Section 2} (timeout jitter) --- same probes, change trial JSON.
    \item \textbf{Section 4 + 7} (audio latency + LED sync) --- move CH1 to LED
          left, add audio jack probe to CH2.
    \item \textbf{Section 6} (valve) --- move CH1 to valve pin.
    \item \textbf{Section 9} (load test) --- repeat 1, 2, 4 with 12 cages running.
    \item \textbf{Sections 3, 5, 10--12} --- secondary, do in a second session.
\end{enumerate}

\end{document}
